\documentclass[11pt,reqno]{amsart}

\usepackage[T1]{fontenc}
\usepackage[utf8]{inputenc}
\usepackage{amsmath,amssymb,amsthm}
\usepackage[letterpaper,margin=0.75in]{geometry}
\usepackage{booktabs}
\usepackage[hidelinks]{hyperref}

\DeclareMathOperator{\vol}{vol}
\DeclareMathOperator{\conv}{conv}
\DeclareMathOperator{\aff}{aff}
\DeclareMathOperator{\area}{area}
\DeclareMathOperator{\spn}{span}
\DeclareMathOperator{\relint}{relint}
\DeclareMathOperator{\interior}{int}

%% Artifact references: a second reference group pointing into the repository
%% rather than the literature, numbered [R1], [R2], ... so that it cannot be
%% confused with the bibliography. Set \repobase and \repocommit once; every
%% link and every citation follows.
\newcommand{\repobase}{https://github.com/csabamedgyeswork/inscribed-polytopes}
\newcommand{\repocommit}{maxvol-d4-n10-v1}
\newcommand{\repohref}{\href{\repobase}{\texttt{\repobase}}}
% \nolinkurl rather than \texttt so a long path may break after a slash.
\newcommand{\artpath}[1]{\href{\repobase/tree/\repocommit/#1}{\nolinkurl{#1}}}

\newcounter{artifact}
\newenvironment{artifactlist}
  {\sloppy % entries are mostly path, so prefer stretched spaces to overrun
   \begin{list}{[R\arabic{artifact}]}
     {\usecounter{artifact}%
      \setlength{\labelwidth}{3.1em}\setlength{\leftmargin}{3.7em}%
      \setlength{\labelsep}{0.6em}\setlength{\itemsep}{0.15ex}%
      \setlength{\parsep}{0pt}}}
  {\end{list}}
\newcommand{\artifact}[3]{\item\label{art:#1}\artpath{#2}, #3}
\newcommand{\artref}[1]{\hyperref[art:#1]{[R\ref*{art:#1}]}}

\newtheorem{theorem}{Theorem}[section]
\newtheorem{proposition}[theorem]{Proposition}
\newtheorem{lemma}[theorem]{Lemma}
\newtheorem{corollary}[theorem]{Corollary}

\theoremstyle{definition}
\newtheorem{definition}[theorem]{Definition}
\newtheorem{remark}[theorem]{Remark}
\newtheorem{hypothesis}[theorem]{Hypothesis}
\newtheorem*{notation}{Notation}

% A quoted statement kept under the name and number of its source.
\newcommand{\citedname}{}
\newtheorem*{citedstatement}{\citedname}
\newenvironment{cited}[2]
  {\renewcommand{\citedname}{#1}\begin{citedstatement}[#2]}
  {\end{citedstatement}}

\title[The inscribed $4$-polytope of maximal volume with ten vertices]
      {The inscribed $4$-polytope of maximal volume\\ with ten vertices}
\author{Csaba Medgyes}
\email{csaba.medgyes.work@gmail.com}

\begin{document}

\begin{abstract}
  Draft for review: definitions and the proof that the method is sound. We
  extend the sum-of-squares certificates of earlier work by nonnegative
  multiples of the facet polynomials $\det(x_{F_{1}} - x_{j}, \ldots,
  x_{F_{4}} - x_{j})$, which are positive on every positively oriented
  realisation of a simplicial type. Such a convexity certificate bounds the
  volume of the realisations of a type without bounding its signed volume on
  all configurations of points on the sphere. We show that if every simplicial
  type with ten vertices other than the join of two $5$-cycles has a
  convexity certificate below $\frac{25}{16}(10 + 2\sqrt{5})$, then the
  inscribed $4$-polytope of maximal volume with ten vertices is the convex
  hull of two regular pentagons inscribed in $\mathbb{S}^{3}$ in orthogonal
  $2$-planes through the origin.
\end{abstract}

\maketitle

\section{Preliminaries}
\label{sec:prelim}

\subsection{Inscribed polytopes}

\begin{cited}{Notation}{{\cite[Section~1]{horvath-langi-2016}}}
  Throughout the paper, by a polytope we mean a convex polytope. The vertex set
  of a polytope $P$ is denoted by $V(P)$. We denote the family of
  $d$-dimensional polytopes, with $n$ vertices and inscribed in the unit sphere
  $\mathbb{S}^{d-1}$, by $\mathcal{P}_{d}(n)$. We denote $d$-dimensional
  volume by $\vol_{d}$, and set
  $v_{d}(n) = \max\{\vol_{d}(P) : P \in \mathcal{P}_{d}(n)\}$.
\end{cited}

\begin{remark}[{\cite[Remark~2.1]{medgyes-maxvol-d4}}]
  \label{rem:attained}
  For $n \geq d + 1$, $v_{d}(n)$ is attained by some
  $P \in \mathcal{P}_{d}(n)$.
\end{remark}

\begin{definition}[{\cite[Definition~2.1, Section~2.2 and Proposition~2.16]{ziegler-1995}}]
  A \emph{facet} of a $d$-polytope is a face of dimension $d-1$, and a
  \emph{ridge} is a face of dimension $d-2$. A polytope is \emph{simplicial} if
  every facet is a simplex. Two polytopes are \emph{combinatorially
  equivalent} if their face lattices are isomorphic; an equivalence class is a
  \emph{combinatorial type}.
\end{definition}

Equivalently, two polytopes are combinatorially equivalent if their
facet--vertex incidence matrices differ only by permutations of rows and
columns \cite[Section~15.1.1]{henk-richter-gebert-ziegler-2017}. For
simplicial polytopes, which is the form in which we use the definition, this
means that a bijection between the vertex sets carries the vertex sets of the
facets of one onto those of the other.

\begin{definition}[{\cite[Definition~1]{horvath-langi-2016}}]
  \label{def:propertyz}
  Let $P \in \mathcal{P}_{d}(n)$ be a $d$-polytope with
  $V(P) = \{p_{1}, p_{2}, \ldots, p_{n}\}$. If for each $i$, there is an open
  set $U_{i} \subset \mathbb{S}^{d-1}$ such that $p_{i} \in U_{i}$, and for any
  $q \in U_{i}$, we have
  \begin{equation*}
    \vol_{d}\bigl(\conv((V(P) \setminus \{p_{i}\}) \cup \{q\})\bigr)
      \leq \vol_{d}(P),
  \end{equation*}
  then we say that $P$ satisfies Property Z.
\end{definition}

\begin{remark}[{\cite{berman-hanes-1970}, \cite[Remark~2.4]{medgyes-maxvol-d4}}]
  \label{rem:propertyz}
  Every $P \in \mathcal{P}_{d}(n)$ with $\vol_{d}(P) = v_{d}(n)$ satisfies
  Property Z.
\end{remark}

\begin{lemma}[{\cite[Lemma~1]{horvath-langi-2016}}]
  \label{lem:simplicial}
  Every $P \in \mathcal{P}_{d}(n)$ satisfying Property Z is simplicial.
\end{lemma}

\subsection{Signed volume and certificates}

For $v_{1}, \ldots, v_{d} \in \mathbb{R}^{d}$, $\det(v_{1}, \ldots, v_{d})$
denotes the determinant of the matrix with columns $v_{1}, \ldots, v_{d}$.

\begin{definition}[{\cite[Definition~3.1]{medgyes-maxvol-d4}}]
  \label{def:signed-volume}
  An \emph{orientation} of a finite set is an ordering of its elements up to
  even permutations. Let $\mathcal{C} = (n, \mathcal{F})$, where
  $\mathcal{F}$ is a finite set of ordered $d$-tuples
  $F = (F_{1}, \ldots, F_{d})$ of distinct elements of $\{1, \ldots, n\}$,
  called facets. Deleting $F_{a}$ gives a ridge, with the orientation of
  $(F_{1}, \ldots, \widehat{F_{a}}, \ldots, F_{d})$ if $a$ is even and the
  opposite one if $a$ is odd. The orientation of $\mathcal{C}$ is
  \emph{coherent} if every ridge lies in exactly two facets and these induce
  opposite orientations on it. A point
  $p = (p_{1}, \ldots, p_{n}) \in (\mathbb{R}^{d})^{n}$ \emph{realises}
  $\mathcal{C}$ if $P = \conv\{p_{1}, \ldots, p_{n}\}$ has the $p_{i}$ as its
  vertices and its facets are exactly the simplices
  $\conv\{p_{j} : j \in F\}$, $F \in \mathcal{F}$. The \emph{signed volume}
  of $p$ is
  \begin{equation}
    \label{eq:signed-volume}
    D(p) = \sum_{F \in \mathcal{F}} \det(p_{F_{1}}, \ldots, p_{F_{d}}).
  \end{equation}
\end{definition}

We write $j \in F$ if $j$ is one of the entries of $F$, and $p_{F}$ for the
set $\{p_{F_{1}}, \ldots, p_{F_{d}}\}$.

\begin{lemma}[{\cite[Lemma~3.2]{medgyes-maxvol-d4}}]
  \label{lem:volume}
  If the orientation of $\mathcal{C}$ is coherent and $p$ realises
  $\mathcal{C}$, then
  \[
    d!\,\vol_{d}(\conv\{p_{1}, \ldots, p_{n}\}) = |D(p)|.
  \]
\end{lemma}

\begin{definition}[{\cite[Definition~3.4]{medgyes-maxvol-d4}}]
  \label{def:certificate}
  Let $x_{i} = (x_{i,1}, \ldots, x_{i,d})$ be the coordinates of $p_{i}$, and
  $v(x)$ the vector of all monomials of degree at most $2$ in the $x_{i,a}$. A
  \emph{certificate} for $\mathcal{C}$ and $\beta \in \mathbb{Q}$ is a
  symmetric positive definite matrix $Q$ with rational entries together with
  polynomials $h_{1}, \ldots, h_{n}$ with rational coefficients such that
  \begin{equation}
    \label{eq:certificate}
    \beta - D(x) = v(x)^{\mathsf{T}} Q\, v(x)
      + \sum_{i=1}^{n} h_{i}(x)\bigl(|x_{i}|^{2} - 1\bigr)
  \end{equation}
  holds identically as polynomials in the $x_{i,a}$.
\end{definition}

\begin{proposition}[{\cite[Proposition~3.5]{medgyes-maxvol-d4}}]
  \label{prop:certificate}
  Suppose the orientation of $\mathcal{C}$ is coherent and a certificate for
  $\mathcal{C}$ and $\beta$ exists. Then $|D(p)| \leq \beta$ for every
  $p \in (\mathbb{S}^{d-1})^{n}$, and every polytope inscribed in
  $\mathbb{S}^{d-1}$ that realises $\mathcal{C}$ has volume at most
  $\beta/d!$.
\end{proposition}

The first conclusion of Proposition~\ref{prop:certificate} concerns every
configuration on the sphere, convex or not. A certificate therefore cannot
exist for $\beta$ below the maximum of $|D|$ on
$(\mathbb{S}^{d-1})^{n}$, even if that maximum is attained only by
configurations that realise no polytope of the type.
Section~\ref{sec:convexity} removes this restriction.

\subsection{The join of two pentagons}

We use the following results of \cite{medgyes-maxvol-d4-n9} with
$m = l = 5$. For $m, l \geq 3$, $\mathcal{J}_{m,l}$ denotes the join of an
$m$-cycle and an $l$-cycle \cite[Definition~3.1]{medgyes-maxvol-d4-n9}, a
coherently oriented type with $m + l$ vertices
\cite[Lemma~3.3]{medgyes-maxvol-d4-n9}.

\begin{lemma}[{\cite[Lemma~3.2]{medgyes-maxvol-d4-n9}}]
  \label{lem:freesum}
  Let $m, l \geq 3$, let $u_{1}, u_{2}, u_{3}, u_{4}$ be an orthonormal basis of
  $\mathbb{R}^{4}$, and for $i \in \mathbb{Z}_{m}$, $j \in \mathbb{Z}_{l}$ let
  \[
    a_{i} = \cos\tfrac{2\pi i}{m}\,u_{1} + \sin\tfrac{2\pi i}{m}\,u_{2}, \qquad
    b_{j} = \cos\tfrac{2\pi j}{l}\,u_{3} + \sin\tfrac{2\pi j}{l}\,u_{4}.
  \]
  Then $P = \conv\{a_{i}, b_{j}\} \in \mathcal{P}_{4}(m + l)$, and the point
  $p = (a_{1}, \ldots, a_{m}, b_{1}, \ldots, b_{l})$ realises
  $\mathcal{J}_{m,l}$.
\end{lemma}

\begin{corollary}[{\cite[Corollary~5.2]{medgyes-maxvol-d4-n9}}]
  \label{cor:pair}
  Let $m, l \in \{3, 4, 5\}$. If $P \in \mathcal{P}_{4}(m+l)$ is
  combinatorially equivalent to the polytope of Lemma~\ref{lem:freesum}, then
  \[
    \vol_{4}(P) \leq \frac{1}{24}\, m\sin\frac{2\pi}{m}\; l\sin\frac{2\pi}{l},
  \]
  with equality if and only if $P$ is the polytope of Lemma~\ref{lem:freesum}
  for some orthonormal basis $u_{1}, \ldots, u_{4}$, that is, up to an
  orthogonal transformation, the convex hull of a regular $m$-gon and a
  regular $l$-gon inscribed in $\mathbb{S}^{3}$ in orthogonal $2$-planes
  through the origin.
\end{corollary}

For $m = l = 5$ the bound is $V^{*} = \frac{25}{24}\sin^{2}\frac{2\pi}{5}$.
Since $\sin^{2}\frac{2\pi}{5} = \frac{1}{16}(10 + 2\sqrt{5})$,
\begin{equation}
  \label{eq:vstar}
  24\,V^{*} = \frac{25}{16}\bigl(10 + 2\sqrt{5}\bigr), \qquad
  V^{*} = \frac{25}{384}\bigl(10 + 2\sqrt{5}\bigr) = 0.942196\ldots.
\end{equation}

\section{Convexity certificates}
\label{sec:convexity}

Throughout this section $\mathcal{C} = (n, \mathcal{F})$ is as in
Definition~\ref{def:signed-volume}, with $d$-tuples as facets.

\subsection{Facet polynomials}

\begin{definition}
  \label{def:facet-polynomial}
  For $F \in \mathcal{F}$ and $j \in \{1, \ldots, n\}$ with $j \notin F$, the
  \emph{facet polynomial} is
  \begin{equation}
    \label{eq:facet-polynomial}
    A_{F,j}(x) = \det\bigl(x_{F_{1}} - x_{j}, \ldots, x_{F_{d}} - x_{j}\bigr).
  \end{equation}
\end{definition}

$A_{F,j}$ is a polynomial of degree $d$ in the $x_{i,a}$ with integer
coefficients, and $|A_{F,j}(p)|$ is $d!$ times the volume of the simplex
$\conv(p_{F} \cup \{p_{j}\})$. An odd permutation of the entries of $F$
changes its sign, so $A_{F,j}$ depends on the ordering of $F$ given in
$\mathcal{F}$, up to even permutations, and not only on its vertex set.

\begin{lemma}
  \label{lem:facet-sign}
  Suppose the orientation of $\mathcal{C}$ is coherent, $p$ realises
  $\mathcal{C}$, and $D(p) > 0$. Then $A_{F,j}(p) > 0$ for every
  $F \in \mathcal{F}$ and every $j \notin F$.
\end{lemma}

\begin{proof}
  Let $P = \conv\{p_{1}, \ldots, p_{n}\}$, and for $F \in \mathcal{F}$ and
  $y \in \mathbb{R}^{d}$ put
  \[
    a_{F}(y) = \det\bigl(p_{F_{1}} - y, \ldots, p_{F_{d}} - y\bigr),
  \]
  so that $A_{F,j}(p) = a_{F}(p_{j})$. Since $p$ realises $\mathcal{C}$,
  $F \mapsto \conv p_{F}$ maps $\mathcal{F}$ onto the facets of $P$. It is
  injective: two tuples with the same vertex set would share all their
  ridges, and a ridge of one of them would lie in at least three tuples,
  the two and the other tuple at that ridge, against coherence.

  \emph{Step 1: $D(p) = \sum_{F \in \mathcal{F}} a_{F}(y)$ for every
  $y \in \mathbb{R}^{d}$.} Expand $a_{F}(y)$ by multilinearity, writing each
  column $p_{F_{a}} - y$ as the sum of $p_{F_{a}}$ and $-y$. Terms with $-y$
  in two columns vanish, since two of their columns are equal. The term
  without $y$ is $\det(p_{F_{1}}, \ldots, p_{F_{d}})$. The term with $-y$ in
  column $a$ is, after moving $y$ to the first column past $a - 1$ columns,
  \[
    -(-1)^{a-1}\det\bigl(y, p_{F_{1}}, \ldots, \widehat{p_{F_{a}}}, \ldots,
    p_{F_{d}}\bigr)
    = (-1)^{a}\det\bigl(y, p_{F_{1}}, \ldots, \widehat{p_{F_{a}}}, \ldots,
    p_{F_{d}}\bigr).
  \]
  For an ordered $(d-1)$-tuple $R$ put
  $\phi_{y}(R) = \det(y, p_{R_{1}}, \ldots, p_{R_{d-1}})$. It is unchanged
  by even permutations of $R$ and changes sign under odd ones, so it is a
  function of the oriented ridge that changes sign when the orientation is
  reversed. By Definition~\ref{def:signed-volume}, the ridge obtained from
  $F$ by deleting $F_{a}$ carries the orientation of
  $(F_{1}, \ldots, \widehat{F_{a}}, \ldots, F_{d})$ if $a$ is even and the
  opposite one if $a$ is odd, so the term above is $\phi_{y}$ of that
  oriented ridge. By coherence, every ridge lies in exactly two facets,
  which induce opposite orientations on it, so in $\sum_{F} a_{F}(y)$ these
  terms cancel in pairs. What remains is
  $\sum_{F} \det(p_{F_{1}}, \ldots, p_{F_{d}}) = D(p)$.

  \emph{Step 2: $a_{F}(y) > 0$ for all $F \in \mathcal{F}$ and all
  $y \in \interior P$.} Fix $y \in \interior P$ and put
  $S_{F} = \conv(\{y\} \cup p_{F})$, so that $|a_{F}(y)| = d!\,\vol_{d}(S_{F})$.
  The simplices $S_{F}$ cover $P$: for $z \in P \setminus \{y\}$, the ray from
  $y$ through $z$ leaves $P$ at a unique point $w \in \partial P$ with
  $z \in [y, w]$, and $w$ lies in some facet $\conv p_{F}$, so $z \in S_{F}$.
  Their interiors are pairwise disjoint: a point $z$ of the interior of
  $S_{F}$ is $z = (1 - t)y + tu$ with $t \in (0, 1)$ and $u$ in the relative
  interior of $\conv p_{F}$, so $u$ is the point $w$ where the ray from $y$
  through $z$ leaves $P$; if $z$ were also interior to $S_{G}$, then $w$
  would lie in the relative interiors of the facets $\conv p_{F}$ and
  $\conv p_{G}$, and distinct faces of a polytope have disjoint relative
  interiors, so $F = G$. Hence, using Step~1 and Lemma~\ref{lem:volume},
  \[
    \Bigl|\sum_{F} a_{F}(y)\Bigr| = |D(p)| = d!\,\vol_{d}(P)
    = \sum_{F} d!\,\vol_{d}(S_{F}) = \sum_{F} |a_{F}(y)|.
  \]
  Equality in the triangle inequality means that the nonzero $a_{F}(y)$
  all have the same sign. None is zero: $a_{F}(y) = 0$ would put $y$ in the
  affine hull of $p_{F}$, a supporting hyperplane of $P$, which misses
  $\interior P$. Since their sum $D(p)$ is positive, every $a_{F}(y)$ is
  positive.

  \emph{Step 3: conclusion.} Fix $F \in \mathcal{F}$ and $j \notin F$. By the
  expansion in Step~1, $a_{F}$ is an affine function of $y$. It vanishes at
  each $p_{F_{k}}$, since then column $k$ is zero, and it is not constant, by
  Step~2; so its zero set is the hyperplane $H_{F} = \aff p_{F}$. As
  $\conv p_{F}$ is a facet of $P$, $P$ lies in one closed half-space of
  $H_{F}$ and $P \cap H_{F} = \conv p_{F}$. By Step~2 that half-space is
  $\{a_{F} \geq 0\}$, so $a_{F}(p_{j}) \geq 0$. If $a_{F}(p_{j}) = 0$, then
  $p_{j} \in P \cap H_{F} = \conv p_{F}$, so the vertex $p_{j}$ of $P$
  would be a vertex of the facet $\conv p_{F}$, that is, $p_{j} = p_{F_{k}}$
  for some $k$; this is impossible, since $j \notin F$ and the $n$ vertices
  $p_{i}$ of $P$ are distinct. Hence $A_{F,j}(p) = a_{F}(p_{j}) > 0$.
\end{proof}

Only $A_{F,j}(p) \geq 0$ is used below. The hypothesis $D(p) > 0$ fixes the
sign: replacing $p$ by its mirror image changes the signs of $D$ and of
every $A_{F,j}$.

\subsection{Convexity certificates}

\begin{definition}
  \label{def:convexity-certificate}
  Let $v_{1}(x), \ldots, v_{r}(x)$ be vectors of monomials of degree at most
  $2$ in the $x_{i,a}$. A \emph{convexity certificate} for $\mathcal{C}$ and
  $\beta \in \mathbb{Q}$ consists of symmetric positive definite matrices
  $Q_{1}, \ldots, Q_{r}$ with rational entries, polynomials
  $h_{1}, \ldots, h_{n}$ with rational coefficients, and rationals
  $\lambda_{F,j} \geq 0$ for $F \in \mathcal{F}$ and $j \notin F$, such that
  \begin{equation}
    \label{eq:convexity-certificate}
    \beta - D(x)
    = \sum_{b=1}^{r} v_{b}(x)^{\mathsf{T}} Q_{b}\, v_{b}(x)
      + \sum_{F \in \mathcal{F}} \sum_{j \notin F} \lambda_{F,j}\, A_{F,j}(x)
      + \sum_{i=1}^{n} h_{i}(x)\bigl(|x_{i}|^{2} - 1\bigr)
  \end{equation}
  holds identically as polynomials in the $x_{i,a}$.
\end{definition}

With $r = 1$, $v_{1} = v$ and every $\lambda_{F,j} = 0$, a convexity
certificate is a certificate in the sense of Definition~\ref{def:certificate}.
Several blocks $v_{b}$ allow the block-diagonal Gram matrices used in the
computations \cite[Section~3.4]{medgyes-maxvol-d4}; a block may omit
monomials, and two blocks may share one.

\begin{proposition}
  \label{prop:convexity-certificate}
  Suppose the orientation of $\mathcal{C}$ is coherent and a convexity
  certificate for $\mathcal{C}$ and $\beta$ exists. Then every polytope
  inscribed in $\mathbb{S}^{d-1}$ that realises $\mathcal{C}$ has volume at
  most $\beta/d!$.
\end{proposition}

\begin{proof}
  Let $P$ be such a polytope, so that $P = \conv\{p_{1}, \ldots, p_{n}\}$
  for some $p \in (\mathbb{S}^{d-1})^{n}$ that realises $\mathcal{C}$. By
  Lemma~\ref{lem:volume}, $|D(p)| = d!\,\vol_{d}(P) > 0$. If $D(p) < 0$,
  let $\rho$ be a reflection in a hyperplane through the origin, and
  replace $p$ by $\rho p = (\rho p_{1}, \ldots, \rho p_{n})$. Its points
  lie on $\mathbb{S}^{d-1}$, and it realises $\mathcal{C}$, since $\rho$
  maps $P$ onto $\conv\{\rho p_{1}, \ldots, \rho p_{n}\}$ and faces onto
  faces. Moreover $\vol_{d}(\rho P) = \vol_{d}(P)$ and
  $D(\rho p) = \det(\rho)\,D(p) = -D(p)$. So we may assume $D(p) > 0$,
  and then $D(p) = d!\,\vol_{d}(P)$.

  Evaluate \eqref{eq:convexity-certificate} at $p$. The last sum vanishes
  since $|p_{i}| = 1$. Each $v_{b}(p)^{\mathsf{T}} Q_{b}\, v_{b}(p)$ is
  nonnegative since $Q_{b}$ is positive definite. Each
  $\lambda_{F,j} A_{F,j}(p)$ is nonnegative since $\lambda_{F,j} \geq 0$ and,
  by Lemma~\ref{lem:facet-sign}, $A_{F,j}(p) > 0$. Hence
  $\beta - d!\,\vol_{d}(P) \geq 0$, that is,
  $\vol_{d}(P) \leq \beta/d!$.
\end{proof}

\begin{remark}
  \label{rem:what-it-bounds}
  Let $K_{\mathcal{C}}$ be the set of $p \in (\mathbb{S}^{d-1})^{n}$ with
  $A_{F,j}(p) \geq 0$ for all $F \in \mathcal{F}$ and $j \notin F$.
  Identity \eqref{eq:convexity-certificate} is a sum-of-squares certificate
  for $D \leq \beta$ on the semialgebraic set $K_{\mathcal{C}}$
  \cite[Section~5]{parrilo-2003}. It says nothing about $D$ outside
  $K_{\mathcal{C}}$, nor about $-D$. The proof above needs neither, since
  every realisation has a mirror image with $D > 0$, which lies in
  $K_{\mathcal{C}}$ by Lemma~\ref{lem:facet-sign}. So, unlike a certificate
  of Definition~\ref{def:certificate}, a convexity certificate can exist for
  $\beta$ below the maximum of $|D|$ on $(\mathbb{S}^{d-1})^{n}$, but not
  below the maximum of $D$ on $K_{\mathcal{C}}$. That maximum is at least
  $d!$ times the largest volume of an inscribed realisation, and may be
  larger, since $K_{\mathcal{C}}$ also contains configurations that realise
  no polytope, for instance ones in which two points coincide.
\end{remark}

\subsection{Verification}
\label{sec:verification}

A convexity certificate is finite rational data: the ordered facets of
$\mathcal{C}$, $\beta$, the blocks $v_{b}$ and entries of the $Q_{b}$, the
coefficients of the $h_{i}$, and the triples $(F, j, \lambda_{F,j})$ with
$\lambda_{F,j} \neq 0$. The verifier checks the hypotheses of
Proposition~\ref{prop:convexity-certificate} in exact rational arithmetic:
\begin{enumerate}
  \item \emph{Multipliers.} For every stored triple, $\lambda_{F,j} \geq 0$,
    $F$ is one of the stored facets with its stored ordering, and
    $j \in \{1, \ldots, n\}$ is not an entry of $F$. The ordering matters,
    since reversing it changes the sign of $A_{F,j}$.
  \item \emph{Identity.} $D$ is rebuilt from the stored facets by
    \eqref{eq:signed-volume} and each $A_{F,j}$ by
    \eqref{eq:facet-polynomial}, both sides of
    \eqref{eq:convexity-certificate} are expanded, and every coefficient
    must agree.
  \item \emph{Positive definiteness.} Each $Q_{b}$ is symmetric, and either
    $M = R^{\mathsf{T}} Q_{b} R$ is strictly diagonally dominant with
    positive diagonal for some upper triangular $R$ with nonzero diagonal
    entries, or $Q_{b}$ has positive leading principal minors, computed by
    the fraction-free elimination of Bareiss \cite{bareiss-1968}. In the
    first case $M$ is positive definite
    \cite[Theorem~6.1.10]{horn-johnson-2013}, and hence so is $Q_{b}$, since
    $R$ is invertible; in the second, Sylvester's criterion
    \cite[Theorem~7.2.5]{horn-johnson-2013} applies. The verifier takes for
    $R$ a rounded floating-point inverse Cholesky factor of $Q_{b}$, with
    integer entries, and computes $M$ exactly; the rounding only chooses $R$.
  \item \emph{Orientation and type.} Every ridge lies in exactly two stored
    facets, with opposite induced orientations, so the orientation is
    coherent; and the vertex sets of the stored facets are exactly those of
    the combinatorial type the certificate is for.
\end{enumerate}
Checks 1 to 3 show that the data is a convexity certificate, and check~4
that $\mathcal{C}$ is coherent and is the intended type. Check~4 needs the
facets of the type to be exact: in Section~\ref{sec:reduction} they come
from polytopes with rational vertices, by exact supporting-hyperplane tests
as in \cite[Lemma~6.1]{medgyes-maxvol-d4-n9}.

\section{Reduction to certificates for ten vertices}
\label{sec:reduction}

\begin{theorem}[{\cite[Theorem~25(i) and Table~3]{firsching-2017}}]
  \label{thm:count}
  There are exactly $162004$ combinatorial types of simplicial $4$-polytopes
  with $10$ vertices.
\end{theorem}

Firsching's data \cite{firsching-2017} give a polytope with rational
vertices for each line of a file with $162004$ lines \artref{data}; we label
them \texttt{firsching-10-0} to \texttt{firsching-10-162003}. The proof of
the main result rests on the following two statements, which are finite
computations. The first is the analogue of
\cite[Lemma~6.1]{medgyes-maxvol-d4-n9} for ten vertices; the second is the
elimination of the other types.

\begin{hypothesis}
  \label{hyp:data}
  The polytopes \texttt{firsching-10-0} to \texttt{firsching-10-162003} are
  simplicial and pairwise not combinatorially equivalent, and for some
  $k_{0}$ the polytope \texttt{firsching-10-$k_{0}$} is combinatorially
  equivalent to the polytope of Lemma~\ref{lem:freesum} with $m = l = 5$.
\end{hypothesis}

\begin{hypothesis}
  \label{hyp:certificates}
  For every $k \neq k_{0}$ there are a coherent orientation $\mathcal{C}_{k}$
  of the facets of \texttt{firsching-10-$k$}, that is, a coherently oriented
  $\mathcal{C}_{k} = (10, \mathcal{F}_{k})$ whose facet vertex sets are
  exactly those of \texttt{firsching-10-$k$}, and a convexity certificate for
  $\mathcal{C}_{k}$ and some $\beta_{k} \in \mathbb{Q}$ with
  $\beta_{k} < \frac{25}{16}(10 + 2\sqrt{5})$.
\end{hypothesis}

By Theorem~\ref{thm:count}, Hypothesis~\ref{hyp:data} says that the
$162004$ polytopes represent every combinatorial type of simplicial
$4$-polytope with ten vertices, each exactly once. The polytopes need not be
inscribed: a convexity certificate is a polynomial identity and needs no
realisation, and for a type without an inscribed realisation the conclusion
of Proposition~\ref{prop:convexity-certificate} holds vacuously.
Hypothesis~\ref{hyp:certificates} allows $\lambda_{F,j} = 0$, so a
certificate in the sense of Definition~\ref{def:certificate} also serves.

\begin{remark}
  \label{rem:exact-bound}
  The condition on $\beta_{k}$ is decided exactly. With
  $t = \frac{16}{25}\beta_{k} - 10$, a rational number,
  $\beta_{k} < \frac{25}{16}(10 + 2\sqrt{5})$ if and only if $t < 2\sqrt{5}$,
  that is, if and only if $t < 0$ or $t^{2} < 20$.
\end{remark}

\begin{proposition}
  \label{prop:rivals}
  Assume Hypotheses~\ref{hyp:data} and~\ref{hyp:certificates}, and let
  $P \in \mathcal{P}_{4}(10)$ be simplicial. If $P$ is not combinatorially
  equivalent to \texttt{firsching-10-$k_{0}$}, then $\vol_{4}(P) < V^{*}$,
  with $V^{*}$ as in \eqref{eq:vstar}.
\end{proposition}

\begin{proof}
  By Hypothesis~\ref{hyp:data} and Theorem~\ref{thm:count}, $P$ is
  combinatorially equivalent to \texttt{firsching-10-$k$} for some
  $k \neq k_{0}$. Since $P$ is simplicial, a bijection between the vertex
  sets carries the facet vertex sets of this polytope, which are those of
  $\mathcal{C}_{k}$, onto those of $P$. So the vertices of $P$ can be listed
  as $p = (p_{1}, \ldots, p_{10})$ so that $p$ realises $\mathcal{C}_{k}$, and
  Proposition~\ref{prop:convexity-certificate} and
  Hypothesis~\ref{hyp:certificates} give
  \[
    \vol_{4}(P) \leq \frac{\beta_{k}}{24}
      < \frac{25}{384}\bigl(10 + 2\sqrt{5}\bigr) = V^{*}.
  \]
\end{proof}

\begin{theorem}
  \label{thm:main}
  Assume Hypotheses~\ref{hyp:data} and~\ref{hyp:certificates}. Then
  \[
    v_{4}(10) = \frac{25}{24}\sin^{2}\frac{2\pi}{5}
      = \frac{25}{384}\bigl(10 + 2\sqrt{5}\bigr) = 0.942196\ldots,
  \]
  and every $P \in \mathcal{P}_{4}(10)$ with $\vol_{4}(P) = v_{4}(10)$ is, up to
  an orthogonal transformation, the convex hull of two regular pentagons
  inscribed in $\mathbb{S}^{3}$ in orthogonal $2$-planes through the origin.
\end{theorem}

\begin{proof}
  By Lemma~\ref{lem:freesum} and Corollary~\ref{cor:pair} with
  $m = l = 5$, the pair of pentagons lies in $\mathcal{P}_{4}(10)$ and has
  volume $V^{*}$, so $v_{4}(10) \geq V^{*}$.

  Let $P \in \mathcal{P}_{4}(10)$ with $\vol_{4}(P) = v_{4}(10)$, which exists
  by Remark~\ref{rem:attained}. By Remark~\ref{rem:propertyz} and
  Lemma~\ref{lem:simplicial}, $P$ is simplicial. Since
  $\vol_{4}(P) \geq V^{*}$, Proposition~\ref{prop:rivals} shows that $P$ is
  combinatorially equivalent to \texttt{firsching-10-$k_{0}$}, and so, by
  Hypothesis~\ref{hyp:data}, to the polytope of Lemma~\ref{lem:freesum} with
  $m = l = 5$. Corollary~\ref{cor:pair} gives
  $v_{4}(10) = \vol_{4}(P) \leq V^{*}$, so $v_{4}(10) = V^{*}$, and its
  equality case identifies $P$.
\end{proof}

\begin{remark}
  \label{rem:computations}
  Hypothesis~\ref{hyp:data} is checked as for nine vertices
  \cite[Lemma~6.1]{medgyes-maxvol-d4-n9}: the facets of each polytope are
  computed from its rational vertices by exact supporting-hyperplane tests,
  each $3$-subset of a facet must lie in exactly two facets, and the
  $162004$ facet systems must have pairwise distinct canonical forms under
  relabelling. Hypothesis~\ref{hyp:certificates} is checked by the
  verification of Section~\ref{sec:verification} for each $k \neq k_{0}$,
  together with Remark~\ref{rem:exact-bound}. The certificates are found by
  semidefinite programming and rounded to rationals as in
  \cite[Section~3.4]{medgyes-maxvol-d4}; none of those steps enters the
  proof.
\end{remark}


\section*{AI assistance}

An AI assistant was used during this work, including for mathematical
discussion, drafting and revising the text, and checking the cited statements
against their sources. The proofs do not depend on the assistant: they rest on
the cited results, on the arguments given here, and on the certificates, which
the exact verifier checks.

\section*{Code and data}
\label{sec:artifacts}

{\footnotesize
Everything reported here is available at \repohref, in the version tagged
\href{\repobase/tree/\repocommit}{\texttt{\repocommit}}. The entries below are cited
as [R$k$]; they are artifacts rather than literature and are numbered separately
from the bibliography.

\begin{artifactlist}
  \artifact{data}{data/polytopalsimplicialspheres10\_4.txt}{Firsching's
    realisations of the $162004$ simplicial types with ten vertices, fetched by
    \artpath{scripts/setup\_data.py}}
  \artifact{verifier}{src/inscribed/certify/verify.py}{the exact verifier of
    \cite[Section~3.3]{medgyes-maxvol-d4}, with
    \artpath{src/inscribed/combinatorics/complex.py} checking coherence}
\end{artifactlist}
}

\bibliographystyle{amsplain}
\bibliography{references}

\end{document}
